\section{Research question, scope, and antecedents}\label{tsc:scope}
The repository's continuation
\src{stieltjes-correlation-calculus/continuations/shifted-hurwitz-jets/}
ends with a concrete question about three distinct translations
\cite{tsc:shj}. Its bilinear theory has only one surviving Fourier index;
three factors impose one relation among three nonzero frequencies and leave
a genuine double sum. The proposed target is an all-index coordinate class,
not an unproved assertion that every resulting period reduces to one-variable
Stieltjes constants.

We resolve that target with two interchangeable representations. The first
is a finite coefficient formula in six colored Tornheim functions
(Theorems~\ref{tsc:master} and \ref{tsc:stieltjes-triple}). The second is an
explicit convergent integral formula using only spectral derivatives of
classical polylogarithms and a finite endpoint correction
(Theorem~\ref{tsc:mellin-all}). Neither definition invokes an unspecified
regularization, an integer-relation search, or a numerical limiting value.
The entire continuation of each colored double sum is proved before any
spectral coefficient is extracted.

The reflected subclass behaves differently from the generic triple.
Cotangent partial fractions force a finite reduction at every number of
separated factors. Combining that algebra with a normalized Stieltjes
transform gives Theorem~\ref{tsc:all-depth}, and rational shifts give
Theorem~\ref{tsc:rational}. The distinction is substantive: the generic
three-factor theorem supplies a colored depth-two representation; the
reflected theorem supplies a finite single-function reduction. The latter
does not imply the former has such a reduction.

\subsection{What is inherited and what is developed here}
The Hurwitz Fourier expansion, the shift relation, Lerch's formula at zero,
and the polygamma reflection law are classical \cite{tsc:dlmf-zeta,tsc:dlmf-gamma}.
Espinosa and Moll develop Hurwitz integral and primitive calculus
\cite{tsc:em1,tsc:em2}. Zhao gives finite double-polylogarithm reductions of
integer colored Tornheim sums \cite{tsc:zhao}. The unshifted cubic
log-Gamma integral already has a Tornheim-derivative evaluation due to
Bailey, D.~Borwein, and J.~M.~Borwein \cite{tsc:bbb}; the canonical ProveIt
chapter explicitly credits it \cite{tsc:gamma-chapter}. We do not claim any
of these results as discoveries of this article.

The repository already has an entire periodic Hurwitz family, bilinear
Stieltjes closure, contact corrections, and normalized primitive identities
\cite{tsc:repo,tsc:shj}. We reprove the normalization needed below so that
the trilinear statement is self-contained. The principal addition is the
separated three-factor proof chain: entire colored continuation, the six
frequency sectors, all-index extraction, convergent Mellin-jet coordinates,
and the resulting primitive and contact laws. The all-depth reflected
formula and its explicit evaluations are accompanying consequences, with
no claim of global priority for their individual specializations.

\subsection{Source snapshot and limits of the audit}
The final observed repository head was
\begin{center}\pin.\end{center}
The key shifted-jet source was read at that revision and has blob identifier
\src{990369336bdf2794aae3057c3231eed4ecadbef3}.
Other selected manuscript reads and the incoming-archive inventory are
recorded in \src{notes/source_snapshot.json}; not every earlier read was
pinned to the final head. A Library search also recovered the same
three-shift question and the abstract of the later nested-harmonic report.
The seven incoming ZIP files were inventoried, not unpacked or analytically
audited in full. This is a targeted continuation, not an audit of every
archived claim.

The current manuscript records a proof of $S_4$, while $S_6$ and the revised
$S_8$ candidate remain conjectural \cite{tsc:repo}. The present argument does
not change those statuses. It also does not turn a representation theorem
into a proof that its coordinates are arithmetically independent.
